\documentclass[11pt,a4paper]{article}
\usepackage[utf8]{inputenc}
\usepackage[T1]{fontenc}
\usepackage[french]{babel}
\usepackage{amsmath,amssymb,amsfonts,amsthm,mathrsfs,mathtools}
\usepackage[top=2cm,bottom=2cm,left=2cm,right=2cm]{geometry}
\usepackage[dvipsnames]{xcolor}
\usepackage{tcolorbox}
\tcbuselibrary{skins,breakable}
\usepackage{tikz}
\usetikzlibrary{calc,shapes,arrows.meta}
\usepackage{enumitem}
\usepackage{fancyhdr}
\usepackage{lastpage}
\usepackage{hyperref}

% ----------------------------------------------------
% Couleurs officielles du projet (Charte LLG / MPSI)
% ----------------------------------------------------
\definecolor{llgblue}{RGB}{20, 50, 115}
\definecolor{llglight}{RGB}{246, 249, 255}

\definecolor{emeraldframe}{RGB}{18, 110, 75}
\definecolor{emeraldback}{RGB}{242, 250, 246}
\definecolor{emeraldtitle}{RGB}{12, 80, 52}

\definecolor{prepaframe}{RGB}{30, 70, 135}
\definecolor{prepaback}{RGB}{246, 249, 255}
\definecolor{prepatitle}{RGB}{20, 50, 105}

\definecolor{alertframe}{RGB}{185, 75, 20}
\definecolor{alertback}{RGB}{255, 249, 245}
\definecolor{alerttitle}{RGB}{145, 55, 12}

% Configuration des hyperliens
\hypersetup{
    colorlinks=true,
    linkcolor=llgblue,
    urlcolor=llgblue,
    citecolor=llgblue
}

% En-têtes et pieds de page
\pagestyle{fancy}
\fancyhf{}
\renewcommand{\headrulewidth}{0.5pt}
\renewcommand{\footrulewidth}{0.4pt}
\fancyhead[L]{\small\textbf{Lycée Louis-le-Grand} -- 1\textsuperscript{ère} Spécialité Mathématiques (1E3)}
\fancyhead[C]{\small\textbf{DST : Probabilités conditionnelles \& Variables aléatoires}}
\fancyhead[R]{\small Jeudi 19 mai 2022}
\fancyfoot[L]{\footnotesize LLG 2021/2022 -- Devoir Sur Table (1E3)}
\fancyfoot[C]{\footnotesize Énoncé \& Corrigé Intégral Rigoureux}
\fancyfoot[R]{\footnotesize Page \thepage\ / \pageref{LastPage}}

% ----------------------------------------------------
% Boîtes tcolorbox personnalisées
% ----------------------------------------------------
\newtcolorbox{correctionbox}[1][]{
  colback=emeraldback,
  colframe=emeraldframe,
  fonttitle=\bfseries\color{white},
  coltitle=white,
  colbacktitle=emeraldtitle,
  title={Corrigé officiel détaillé et barème indicatif #1},
  arc=2mm,
  boxrule=0.8pt,
  titlerule=0.5pt,
  breakable,
  before skip=0.4cm,
  after skip=0.5cm
}

\newtcolorbox{remarquebox}[1][]{
  colback=prepaback,
  colframe=prepaframe,
  fonttitle=\bfseries\small\color{white},
  coltitle=white,
  colbacktitle=prepatitle,
  title={Prolongement \& Regard Prépa / MPSI #1},
  arc=1.5mm,
  boxrule=0.6pt,
  titlerule=0.4pt,
  breakable,
  before skip=0.25cm,
  after skip=0.35cm
}

\newtcolorbox{alertebox}[1][]{
  colback=alertback,
  colframe=alertframe,
  fonttitle=\bfseries\small\color{white},
  coltitle=white,
  colbacktitle=alerttitle,
  title={Point de vigilance méthodologique \& Pièges fréquents #1},
  arc=1.5mm,
  boxrule=0.6pt,
  titlerule=0.4pt,
  breakable,
  before skip=0.25cm,
  after skip=0.35cm
}

% Commandes mathématiques usuelles
\newcommand{\N}{\mathbb{N}}
\newcommand{\Z}{\mathbb{Z}}
\newcommand{\Q}{\mathbb{Q}}
\newcommand{\R}{\mathbb{R}}
\newcommand{\C}{\mathbb{C}}
\newcommand{\Prob}{\mathbb{P}}
\newcommand{\E}{\mathbb{E}}
\newcommand{\V}{\mathbb{V}}
\newcommand{\dd}{\mathrm{d}}
\newcommand{\e}{\mathrm{e}}
\newcommand{\un}{\mathbf{1}}

\begin{document}

% ====================================================
% EN-TÊTE OFFICIEL DU DEVOIR
% ====================================================
\begin{center}
    {\LARGE \bfseries Lycée Louis-le-Grand -- Paris}\\[0.2cm]
    {\large Classe de Première 3 (1E3) -- Année scolaire 2021/2022}\\[0.3cm]
    \begin{tcolorbox}[colback=llglight,colframe=llgblue,arc=3mm,boxrule=1.2pt,center]
        \begin{center}
            {\Large \bfseries \color{llgblue} DEVOIR SUR TABLE DE MATHÉMATIQUES}\\[0.15cm]
            {\large \bfseries Probabilités conditionnelles, Indépendance, Arbres \& Variables aléatoires}\\[0.2cm]
            \small Date : Jeudi 19 mai 2022 \quad $\vert$ \quad Durée : 2 heures \quad $\vert$ \quad Calculatrice autorisée selon la réglementation
        \end{center}
    \end{tcolorbox}
\end{center}

\vspace{0.3cm}

\begin{tcolorbox}[colback=white,colframe=llgblue!60,arc=1.5mm,boxrule=0.6pt]
\small\textbf{Consignes générales d'évaluation :}
\begin{itemize}[leftmargin=*,itemsep=1pt]
    \item La clarté des raisonnements, la précision de la rédaction et la rigueur dans l'emploi du vocabulaire ensembliste et probabiliste constituent des critères déterminants d'appréciation.
    \item Tous les arbres de probabilités doivent être construits avec soin, en faisant figurer les événements sur les nœuds et les probabilités conditionnelles exactes sur chaque branche.
    \item Toute formule générale de cours (formule des probabilités totales, formule de Bayes, définition de l'espérance mathématique) doit être explicitement citée avant son application numérique.
\end{itemize}
\end{tcolorbox}

\vspace{0.4cm}

% ====================================================
% PARTIE I : ÉNONCÉ INTÉGRAL ÉLÈVE (BARÈME /20)
% ====================================================
\begin{center}
    {\Large \bfseries \color{llgblue} \underline{Sujet du Devoir Surveillé (Barème sur 20 points)}}
\end{center}

\vspace{0.2cm}

% ----------------------------------------------------
% Exercice 1 (5 points)
% ----------------------------------------------------
\subsection*{Exercice 1 \hfill (5 points)}

Le directeur d'une grande entreprise a proposé à l'ensemble de ses salariés un stage de formation à l'utilisation d'un nouveau logiciel. Ce stage a été suivi par $25\,\%$ des salariés.

Dans cette entreprise, $52\,\%$ des salariés sont des femmes, parmi lesquelles $40\,\%$ ont suivi le stage.

On interroge au hasard un salarié de l'entreprise et on considère les événements :
\begin{itemize}[label=$\bullet$]
    \item $F$ : \og le salarié interrogé est une femme \fg{},
    \item $S$ : \og le salarié interrogé a suivi le stage \fg{}.
\end{itemize}
$\overline{F}$ et $\overline{S}$ désignent respectivement les événements contraires des événements $F$ et $S$.

\begin{enumerate}
    \item Donner la probabilité de l'événement $S$.
    \item Recopier et compléter les pointillés de l'arbre pondéré ci-dessous sur les quatre branches indiquées :
    \begin{center}
    \begin{tikzpicture}[xscale=2.4, yscale=1.1, thick, >=Stealth]
        \node (R) at (0, 0) {};
        \node (F) at (1.5, 0.8) {$F$};
        \node (Fbar) at (1.5, -0.8) {$\overline{F}$};
        \node (FS) at (3.2, 1.2) {$S$};
        \node (FSbar) at (3.2, 0.4) {$\overline{S}$};
        \node (FbarS) at (3.2, -0.4) {$S$};
        \node (FbarSbar) at (3.2, -1.2) {$\overline{S}$};
        
        \draw (R) -- (F) node[midway, above left] {$\dots$};
        \draw (R) -- (Fbar) node[midway, below left] {$\dots$};
        \draw (F) -- (FS) node[midway, above] {$\dots$};
        \draw (F) -- (FSbar) node[midway, below] {$\dots$};
        \draw (Fbar) -- (FbarS) node[midway, above] {$\dots$};
        \draw (Fbar) -- (FbarSbar) node[midway, below] {$\dots$};
    \end{tikzpicture}
    \end{center}
    \item Démontrer que la probabilité que la personne interrogée soit une femme ayant suivi le stage est égale à $0{,}208$.
    \item On sait que la personne interrogée a suivi le stage. Quelle est la probabilité que ce soit une femme ?
    \item Le directeur affirme que, parmi les hommes salariés de l'entreprise, moins de $10\,\%$ ont suivi le stage. Justifier l'affirmation du directeur.
\end{enumerate}

\vspace{0.3cm}

% ----------------------------------------------------
% Exercice 2 (4 points)
% ----------------------------------------------------
\subsection*{Exercice 2 \hfill (4 points)}

Au cours de la fabrication d'une paire de lunettes, la paire de verres doit subir deux traitements notés $\mathrm{T}_1$ et $\mathrm{T}_2$.

On prélève au hasard une paire de verres dans la production. On désigne par :
\begin{itemize}[label=$\bullet$]
    \item $A$ l'événement : \og la paire de verres présente un défaut pour le traitement $\mathrm{T}_1$ \fg{},
    \item $B$ l'événement : \og la paire de verres présente un défaut pour le traitement $\mathrm{T}_2$ \fg{}.
\end{itemize}
On note respectivement $\overline{A}$ et $\overline{B}$ les événements contraires de $A$ et $B$.

Une étude a montré que :
\begin{itemize}[label=$\bullet$]
    \item la probabilité qu'une paire de verres présente un défaut pour le traitement $\mathrm{T}_1$, notée $P(A)$, est égale à $0{,}1$.
    \item la probabilité qu'une paire de verres présente un défaut pour le traitement $\mathrm{T}_2$, notée $P(B)$, est égale à $0{,}2$.
    \item la probabilité qu'une paire de verres ne présente aucun des deux défauts est $0{,}75$.
\end{itemize}

\begin{enumerate}
    \item Recopier et compléter le tableau suivant avec les probabilités correspondantes :
    \begin{center}
    \renewcommand{\arraystretch}{1.3}
    \begin{tabular}{|c|c|c|c|}
        \hline
        & $A$ & $\overline{A}$ & Total \\
        \hline
        $B$ & $\dots$ & $\dots$ & $\dots$ \\
        \hline
        $\overline{B}$ & $\dots$ & $\dots$ & $\dots$ \\
        \hline
        Total & $\dots$ & $\dots$ & $1$ \\
        \hline
    \end{tabular}
    \end{center}
    \item 
    \begin{enumerate}
        \item Déterminer, en justifiant la réponse, la probabilité qu'une paire de verres, prélevée au hasard dans la production, présente un défaut pour au moins un des deux traitements $\mathrm{T}_1$ ou $\mathrm{T}_2$.
        \item Donner la probabilité qu'une paire de verres, prélevée au hasard dans la production, présente deux défauts, un pour chaque traitement $\mathrm{T}_1$ et $\mathrm{T}_2$.
        \item Les événements $A$ et $B$ sont-ils indépendants ? Justifier la réponse.
    \end{enumerate}
    \item Calculer la probabilité qu'une paire de verres, prélevée au hasard dans la production, présente un défaut pour un seul des deux traitements.
    \item Calculer la probabilité qu'une paire de verres, prélevée au hasard dans la production, présente un défaut pour le traitement $\mathrm{T}_2$, sachant que cette paire de verres présente un défaut pour le traitement $\mathrm{T}_1$.
\end{enumerate}

\vspace{0.3cm}

% ----------------------------------------------------
% Exercice 3 (4,5 points)
% ----------------------------------------------------
\subsection*{Exercice 3 \hfill (4,5 points)}

Le coyote est un animal sauvage proche du loup, qui vit en Amérique du Nord. Dans l'état d'Oklahoma, aux États-Unis, $70\,\%$ des coyotes sont touchés par une maladie appelée \emph{ehrlichiose}.

Il existe un test aidant à la détection de cette maladie. Lorsque ce test est appliqué à un coyote, son résultat est soit positif, soit négatif, et on sait que :
\begin{itemize}[label=$\bullet$]
    \item Si le coyote est malade, le test est positif dans $97\,\%$ des cas.
    \item Si le coyote n'est pas malade, le test est négatif dans $95\,\%$ des cas.
\end{itemize}

\textbf{Partie A}

Des vétérinaires capturent un coyote d'Oklahoma au hasard et lui font subir un test pour l'ehrlichiose. On considère les événements suivants :
\begin{itemize}[label=$\bullet$]
    \item $M$ : \og le coyote est malade \fg{},
    \item $T$ : \og le test du coyote est positif \fg{}.
\end{itemize}
On note $\overline{M}$ et $\overline{T}$ respectivement les événements contraires de $M$ et $T$.

\begin{enumerate}
    \item Recopier et compléter l'arbre pondéré ci-dessous qui modélise la situation :
    \begin{center}
    \begin{tikzpicture}[xscale=2.4, yscale=1.1, thick, >=Stealth]
        \node (R) at (0, 0) {};
        \node (M) at (1.5, 0.8) {$M$};
        \node (Mbar) at (1.5, -0.8) {$\overline{M}$};
        \node (MT) at (3.2, 1.2) {$T$};
        \node (MTbar) at (3.2, 0.4) {$\overline{T}$};
        \node (MbarT) at (3.2, -0.4) {$T$};
        \node (MbarTbar) at (3.2, -1.2) {$\overline{T}$};
        
        \draw (R) -- (M) node[midway, above left] {$\dots$};
        \draw (R) -- (Mbar) node[midway, below left] {$\dots$};
        \draw (M) -- (MT) node[midway, above] {$\dots$};
        \draw (M) -- (MTbar) node[midway, below] {$\dots$};
        \draw (Mbar) -- (MbarT) node[midway, above] {$\dots$};
        \draw (Mbar) -- (MbarTbar) node[midway, below] {$\dots$};
    \end{tikzpicture}
    \end{center}
    \item Déterminer la probabilité que le coyote soit malade et que son test soit positif.
    \item Démontrer que la probabilité de $T$ est égale à $0{,}694$.
    \item On appelle \og valeur prédictive positive du test \fg{} la probabilité que le coyote soit effectivement malade sachant que son test est positif. Calculer la valeur prédictive positive du test. On arrondira le résultat au millième.
    \item 
    \begin{enumerate}
        \item Par analogie avec la question précédente, proposer une définition de la \og valeur prédictive négative du test \fg{} et calculer cette valeur en arrondissant au millième.
        \item Comparer les valeurs prédictives positive et négative du test, et interpréter.
    \end{enumerate}
\end{enumerate}

\vspace{0.3cm}

% ----------------------------------------------------
% Exercice 4 (5 points)
% ----------------------------------------------------
\subsection*{Exercice 4 \hfill (5 points)}

Une urne contient des jetons blancs et noirs tous indiscernables au toucher. Une partie consiste à prélever au hasard successivement et avec remise deux jetons de cette urne.

On établit la règle de jeu suivante :
\begin{itemize}[label=$\bullet$]
    \item un joueur perd $9$ euros si les deux jetons tirés sont de couleur blanche ;
    \item un joueur perd $1$ euro si les deux jetons tirés sont de couleur noire ;
    \item un joueur gagne $5$ euros si les deux jetons tirés sont de couleurs différentes.
\end{itemize}

\begin{enumerate}
    \item On considère que l'urne contient $2$ jetons noirs et $3$ jetons blancs.
    \begin{enumerate}
        \item Modéliser la situation à l'aide d'un arbre pondéré.
        \item Calculer la probabilité de perdre $9$ \euro{} sur une partie.
    \end{enumerate}
    \item On considère maintenant que l'urne contient $3$ jetons blancs et au moins deux jetons noirs mais on ne connaît pas le nombre exact de jetons noirs. On appellera $N$ le nombre de jetons noirs ($N \ge 2$).
    \begin{enumerate}
        \item Soit $X$ la variable aléatoire donnant le gain du jeu pour une partie. Déterminer la loi de probabilité de cette variable aléatoire.
        \item Résoudre l'inéquation pour $x$ réel :
        \[
        -x^2 + 30x - 81 > 0
        \]
        \item En utilisant le résultat de la question précédente, déterminer le nombre de jetons noirs que l'urne doit contenir afin que ce jeu soit favorable au joueur.
        \item Combien de jetons noirs le joueur doit-il demander afin d'obtenir un gain moyen maximal ?
    \end{enumerate}
\end{enumerate}

\vspace{0.3cm}

% ----------------------------------------------------
% Exercice 5 (1,5 point)
% ----------------------------------------------------
\subsection*{Exercice 5 \hfill (1,5 point)}

$A$ et $B$ sont deux avions ayant respectivement $2$ et $4$ moteurs. Les moteurs sont supposés indépendants les uns des autres, et ils ont une probabilité $p \in [0, 1]$ de tomber en panne.

Chaque avion arrive à destination si \textbf{moins de la moitié} de ses moteurs tombe en panne.

Quel avion choisissez-vous ? (On discutera en fonction de $p$).

\newpage

% ====================================================
% PARTIE II : CORRIGÉ OFFICIEL DÉTAILLÉ & PROLONGEMENTS MPSI
% ====================================================
\begin{center}
    {\Large \bfseries \color{emeraldtitle} \underline{Corrigé Intégral Officiel \& Approfondissements Prépa}}
\end{center}

\vspace{0.3cm}

% ----------------------------------------------------
% Correction Exercice 1
% ----------------------------------------------------
\subsection*{Correction détaillée de l'Exercice 1 \hfill (5 points)}

\begin{correctionbox}[-- Analyse ensembliste, arbre pondéré \& Formule de Bayes]
\textbf{1. Probabilité de l'événement $S$ :}
D'après l'énoncé, le stage de formation proposé par la direction a été suivi par $25\,\%$ de l'ensemble des salariés de l'entreprise. Ainsi :
\[
\mathbf{P(S) = 0{,}25}
\]

\textbf{2. Arbre pondéré modélisant la situation :}
Identifions précisément les probabilités données par le texte :
\begin{itemize}[label=$\bullet$]
    \item La proportion de femmes dans l'entreprise est de $52\,\%$, d'où $P(F) = 0{,}52$. Par la loi de l'événement contraire :
    \[
    P(\overline{F}) = 1 - P(F) = 1 - 0{,}52 = 0{,}48
    \]
    \item Parmi les femmes, $40\,\%$ ont suivi le stage, ce qui traduit la probabilité conditionnelle :
    \[
    P_F(S) = 0{,}40 \implies P_F(\overline{S}) = 1 - P_F(S) = 1 - 0{,}40 = 0{,}60
    \]
    \item Pour les branches issues du nœud $\overline{F}$, la probabilité conditionnelle $P_{\overline{F}}(S)$ s'obtient à l'aide de la formule des probabilités totales appliquée à la partition $\{F, \overline{F}\}$ :
    \[
    P(S) = P(F \cap S) + P(\overline{F} \cap S) = P(F) \times P_F(S) + P(\overline{F}) \times P_{\overline{F}}(S)
    \]
    En injectant les valeurs connues :
    \[
    0{,}25 = 0{,}52 \times 0{,}40 + 0{,}48 \times P_{\overline{F}}(S) = 0{,}208 + 0{,}48 \times P_{\overline{F}}(S)
    \]
    On en déduit :
    \[
    0{,}48 \times P_{\overline{F}}(S) = 0{,}25 - 0{,}208 = 0{,}042 \implies P_{\overline{F}}(S) = \frac{0{,}042}{0{,}48} = \frac{42}{480} = \frac{7}{80} = \mathbf{0{,}0875}
    \]
    D'où son contraire :
    \[
    P_{\overline{F}}(\overline{S}) = 1 - 0{,}0875 = \mathbf{0{,}9125}
    \]
\end{itemize}

L'arbre pondéré complet s'établit donc ainsi :
\begin{center}
\begin{tikzpicture}[xscale=2.6, yscale=1.2, thick, >=Stealth]
    \node[circle, fill=llgblue!20, inner sep=2pt] (R) at (0, 0) {$\Omega$};
    \node (F) at (1.5, 0.9) {$F$};
    \node (Fbar) at (1.5, -0.9) {$\overline{F}$};
    \node (FS) at (3.3, 1.4) {$S$};
    \node (FSbar) at (3.3, 0.5) {$\overline{S}$};
    \node (FbarS) at (3.3, -0.5) {$S$};
    \node (FbarSbar) at (3.3, -1.4) {$\overline{S}$};
    
    \draw[->] (R) -- (F) node[midway, above left, text=llgblue] {$0{,}52$};
    \draw[->] (R) -- (Fbar) node[midway, below left, text=llgblue] {$0{,}48$};
    
    \draw[->] (F) -- (FS) node[midway, above, text=emeraldframe] {$0{,}40$};
    \draw[->] (F) -- (FSbar) node[midway, below, text=emeraldframe] {$0{,}60$};
    
    \draw[->] (Fbar) -- (FbarS) node[midway, above, text=alertframe] {$0{,}0875$};
    \draw[->] (Fbar) -- (FbarSbar) node[midway, below, text=alertframe] {$0{,}9125$};
\end{tikzpicture}
\end{center}

\textbf{3. Calcul de $P(F \cap S)$ :}
D'après la formule des probabilités composées :
\[
P(F \cap S) = P(F) \times P_F(S) = 0{,}52 \times 0{,}40 = \mathbf{0{,}208}
\]
La probabilité que le salarié interrogé soit une femme ayant suivi le stage est exactement égale à $0{,}208$ (soit $20{,}8\,\%$ des salariés).

\textbf{4. Probabilité conditionnelle inverse $P_S(F)$ (Formule de Bayes) :}
Sachant que le salarié a suivi le stage, la probabilité qu'il s'agisse d'une femme est donnée par la définition de la probabilité conditionnelle :
\[
P_S(F) = \frac{P(F \cap S)}{P(S)} = \frac{0{,}208}{0{,}25} = \frac{208}{250} = \frac{104}{125} = \mathbf{0{,}832}
\]
Ainsi, $83{,}2\,\%$ des personnes ayant suivi le stage sont des femmes.

\textbf{5. Justification de l'affirmation du directeur :}
Le directeur affirme que parmi les hommes salariés de l'entreprise, moins de $10\,\%$ ont suivi le stage.
L'ensemble des hommes salariés correspond à l'événement $\overline{F}$. La proportion d'hommes ayant suivi le stage est donnée par la probabilité conditionnelle $P_{\overline{F}}(S)$.
Par le calcul mené à la question 2 :
\[
P_{\overline{F}}(S) = \frac{P(\overline{F} \cap S)}{P(\overline{F})} = \frac{P(S) - P(F \cap S)}{P(\overline{F})} = \frac{0{,}25 - 0{,}208}{0{,}48} = \frac{0{,}042}{0{,}48} = 0{,}0875 = 8{,}75\,\%
\]
Comme $8{,}75\,\% < 10\,\%$, \textbf{l'affirmation du directeur est parfaitement justifiée}.
\end{correctionbox}

\begin{remarquebox}[-- Formalisme de Kolmogorov \& Théorème de Bayes]
Soit $(\Omega, \mathscr{T}, P)$ un espace probabilisé. Si $(A_i)_{i \in I}$ forme un système complet d'événements (partition dénombrable de $\Omega$ par des événements de probabilités non nulles), alors pour tout événement $B$ tel que $P(B) > 0$ et tout indice $k \in I$ :
\[
P(A_k \mid B) = \frac{P(A_k \cap B)}{P(B)} = \frac{P(A_k) P(B \mid A_k)}{\displaystyle\sum_{i \in I} P(A_i) P(B \mid A_i)}
\]
Cette formule constitue l'\emph{inversion de Bayes}. En pratique, elle permet d'actualiser la distribution \emph{a priori} $(P(A_i))$ en une distribution \emph{a posteriori} $(P(A_i \mid B))$ dès lors qu'un symptôme ou résultat $B$ est observé. Ici, le fait de savoir qu'un salarié a suivi le stage fait passer la probabilité qu'il soit une femme de $P(F) = 0{,}52$ à $P_S(F) = 0{,}832$.
\end{remarquebox}

\vspace{0.4cm}

% ----------------------------------------------------
% Correction Exercice 2
% ----------------------------------------------------
\subsection*{Correction détaillée de l'Exercice 2 \hfill (4 points)}

\begin{correctionbox}[-- Tableau de contingence, indépendance \& corrélation d'événements]
\textbf{1. Tableau de contingence complété :}
Rappelons les données de l'énoncé :
\begin{itemize}[label=$\bullet$]
    \item $P(A) = 0{,}1 \implies P(\overline{A}) = 1 - 0{,}1 = 0{,}9$.
    \item $P(B) = 0{,}2 \implies P(\overline{B}) = 1 - 0{,}2 = 0{,}8$.
    \item \og Aucun des deux défauts \fg{} correspond à l'événement $\overline{A} \cap \overline{B}$, d'où $P(\overline{A} \cap \overline{B}) = 0{,}75$.
\end{itemize}
En complétant les lignes et colonnes par soustraction :
\begin{itemize}[label=$\bullet$]
    \item Sur la ligne $\overline{B}$ : $P(A \cap \overline{B}) = P(\overline{B}) - P(\overline{A} \cap \overline{B}) = 0{,}80 - 0{,}75 = \mathbf{0{,}05}$.
    \item Sur la colonne $A$ : $P(A \cap B) = P(A) - P(A \cap \overline{B}) = 0{,}10 - 0{,}05 = \mathbf{0{,}05}$.
    \item Sur la ligne $B$ : $P(\overline{A} \cap B) = P(B) - P(A \cap B) = 0{,}20 - 0{,}05 = \mathbf{0{,}15}$.
    \item Vérification sur la colonne $\overline{A}$ : $P(\overline{A} \cap B) + P(\overline{A} \cap \overline{B}) = 0{,}15 + 0{,}75 = 0{,}90 = P(\overline{A})$.
\end{itemize}

Le tableau s'écrit de manière rigoureuse :
\begin{center}
\renewcommand{\arraystretch}{1.35}
\begin{tabular}{|c|c|c|c|}
    \hline
    & \textbf{$A$} & \textbf{$\overline{A}$} & \textbf{Total} \\
    \hline
    \textbf{$B$} & $\mathbf{0{,}05}$ & $\mathbf{0{,}15}$ & $\mathbf{0{,}20}$ \\
    \hline
    \textbf{$\overline{B}$} & $\mathbf{0{,}05}$ & $\mathbf{0{,}75}$ & $\mathbf{0{,}80}$ \\
    \hline
    \textbf{Total} & $\mathbf{0{,}10}$ & $\mathbf{0{,}90}$ & $\mathbf{1{,}00}$ \\
    \hline
\end{tabular}
\end{center}

\textbf{2. (a) Probabilité de présenter au moins un défaut :}
L'événement \og présenter au moins un défaut \fg{} est la réunion $A \cup B$. Son événement contraire est \og ne présenter aucun défaut \fg{}, soit $\overline{A} \cap \overline{B}$. Par les lois de De Morgan :
\[
\overline{A \cup B} = \overline{A} \cap \overline{B}
\]
D'où :
\[
P(A \cup B) = 1 - P(\overline{A} \cap \overline{B}) = 1 - 0{,}75 = \mathbf{0{,}25}
\]
\emph{Remarque de contrôle :} Par la formule de Poincaré :
\[
P(A \cup B) = P(A) + P(B) - P(A \cap B) = 0{,}10 + 0{,}20 - 0{,}05 = 0{,}25
\]

\textbf{2. (b) Probabilité de présenter deux défauts :}
L'événement \og présenter deux défauts, un pour chaque traitement \fg{} est l'intersection $A \cap B$.
D'après la lecture directe du tableau de contingence :
\[
\mathbf{P(A \cap B) = 0{,}05}
\]

\textbf{2. (c) Test d'indépendance de $A$ et $B$ :}
Deux événements $A$ et $B$ sont indépendants si et seulement si $P(A \cap B) = P(A) \times P(B)$.
Calculons le produit des probabilités marginales :
\[
P(A) \times P(B) = 0{,}10 \times 0{,}20 = 0{,}02
\]
Or, nous avons établi que :
\[
P(A \cap B) = 0{,}05
\]
Comme $P(A \cap B) \neq P(A) \times P(B)$ ($0{,}05 \neq 0{,}02$), \textbf{les événements $A$ et $B$ ne sont pas indépendants}.

\textbf{3. Probabilité de présenter un défaut pour un seul des deux traitements :}
Cet événement correspond à la différence symétrique $A \mathbin{\Delta} B = (A \cap \overline{B}) \cup (\overline{A} \cap B)$.
Puisque ces deux événements sont incompatibles :
\[
P(A \mathbin{\Delta} B) = P(A \cap \overline{B}) + P(\overline{A} \cap B) = 0{,}05 + 0{,}15 = \mathbf{0{,}20}
\]
On peut également l'obtenir par : $P(A \cup B) - P(A \cap B) = 0{,}25 - 0{,}05 = 0{,}20$.

\textbf{4. Probabilité conditionnelle $P_A(B)$ :}
Il s'agit de la probabilité que la paire de verres présente le défaut $\mathrm{T}_2$ sachant qu'elle présente le défaut $\mathrm{T}_1$ :
\[
P_A(B) = \frac{P(A \cap B)}{P(A)} = \frac{0{,}05}{0{,}10} = \mathbf{0{,}50}
\]
\end{correctionbox}

\begin{remarquebox}[-- Covariance d'indicatrices et corrélation positive]
Notons $\un_A$ et $\un_B$ les variables aléatoires indicatrices des événements $A$ et $B$.
Leur covariance s'écrit :
\[
\operatorname{Cov}(\un_A, \un_B) = \E(\un_A \un_B) - \E(\un_A)\E(\un_B) = P(A \cap B) - P(A)P(B) = 0{,}05 - 0{,}02 = +0{,}03 > 0
\]
Le signe strictement positif indique une \emph{corrélation positive} entre les défauts : la présence d'un défaut $\mathrm{T}_1$ augmente drastiquement la survenue du défaut $\mathrm{T}_2$, faisant passer sa probabilité de $P(B) = 20\,\%$ à $P_A(B) = 50\,\%$. Industriellement, cela révèle qu'une anomalie en amont sur la ligne de traitement fragilise le verre pour le traitement suivant.
\end{remarquebox}

\vspace{0.4cm}

% ----------------------------------------------------
% Correction Exercice 3
% ----------------------------------------------------
\subsection*{Correction détaillée de l'Exercice 3 \hfill (4,5 points)}

\begin{correctionbox}[-- Dépistage épidémiologique : Sensibilité, Spécificité \& VPP/VPN]
\textbf{1. Arbre pondéré complet modélisant le test :}
Identifions les paramètres épidémiologiques :
\begin{itemize}[label=$\bullet$]
    \item Prévalence de l'ehrlichiose : $P(M) = 0{,}70 \implies P(\overline{M}) = 1 - 0{,}70 = 0{,}30$.
    \item Sensibilité du test : $P_M(T) = 0{,}97 \implies P_M(\overline{T}) = 1 - 0{,}97 = 0{,}03$ (taux de faux négatifs parmi les malades).
    \item Spécificité du test : $P_{\overline{M}}(\overline{T}) = 0{,}95 \implies P_{\overline{M}}(T) = 1 - 0{,}95 = 0{,}05$ (taux de faux positifs parmi les animaux sains).
\end{itemize}

L'arbre pondéré s'établit ainsi :
\begin{center}
\begin{tikzpicture}[xscale=2.6, yscale=1.2, thick, >=Stealth]
    \node[circle, fill=llgblue!20, inner sep=2pt] (R) at (0, 0) {$\Omega$};
    \node (M) at (1.5, 0.9) {$M$};
    \node (Mbar) at (1.5, -0.9) {$\overline{M}$};
    \node (MT) at (3.3, 1.4) {$T$};
    \node (MTbar) at (3.3, 0.5) {$\overline{T}$};
    \node (MbarT) at (3.3, -0.5) {$T$};
    \node (MbarTbar) at (3.3, -1.4) {$\overline{T}$};
    
    \draw[->] (R) -- (M) node[midway, above left, text=llgblue] {$0{,}70$};
    \draw[->] (R) -- (Mbar) node[midway, below left, text=llgblue] {$0{,}30$};
    
    \draw[->] (M) -- (MT) node[midway, above, text=emeraldframe] {$0{,}97$};
    \draw[->] (M) -- (MTbar) node[midway, below, text=emeraldframe] {$0{,}03$};
    
    \draw[->] (Mbar) -- (MbarT) node[midway, above, text=alertframe] {$0{,}05$};
    \draw[->] (Mbar) -- (MbarTbar) node[midway, below, text=alertframe] {$0{,}95$};
\end{tikzpicture}
\end{center}

\textbf{2. Probabilité de l'événement $M \cap T$ :}
D'après la formule des probabilités composées :
\[
P(M \cap T) = P(M) \times P_M(T) = 0{,}70 \times 0{,}97 = \mathbf{0{,}679}
\]

\textbf{3. Démontrer que $P(T) = 0{,}694$ :}
Les événements $M$ et $\overline{M}$ forment une partition de l'univers $\Omega$. D'après la formule des probabilités totales :
\[
P(T) = P(M \cap T) + P(\overline{M} \cap T) = P(M) \times P_M(T) + P(\overline{M}) \times P_{\overline{M}}(T)
\]
En injectant les valeurs numériques :
\[
P(T) = 0{,}679 + 0{,}30 \times 0{,}05 = 0{,}679 + 0{,}015 = \mathbf{0{,}694}
\]
Ce résultat confirme rigoureusement la valeur attendue.

\textbf{4. Calcul de la Valeur Prédictive Positive (VPP) :}
La valeur prédictive positive est la probabilité conditionnelle $P_T(M)$ :
\[
P_T(M) = \frac{P(M \cap T)}{P(T)} = \frac{0{,}679}{0{,}694} = \frac{679}{694} \approx 0{,}978386\dots
\]
Au millième près :
\[
\mathbf{\mathrm{VPP} = P_T(M) \approx 0{,}978} \quad (\text{soit } 97{,}8\,\%)
\]

\textbf{5. (a) Définition et calcul de la Valeur Prédictive Négative (VPN) :}
\begin{itemize}[label=$\bullet$]
    \item \textbf{Définition :} La valeur prédictive négative du test, notée $\mathrm{VPN}$, est la probabilité que le coyote soit effectivement sain (non malade) sachant que son test est négatif, soit la probabilité conditionnelle $P_{\overline{T}}(\overline{M})$.
    \item \textbf{Calcul :} Par définition :
    \[
    P_{\overline{T}}(\overline{M}) = \frac{P(\overline{M} \cap \overline{T})}{P(\overline{T})} = \frac{P(\overline{M}) \times P_{\overline{M}}(\overline{T})}{1 - P(T)}
    \]
    Avec $P(\overline{T}) = 1 - 0{,}694 = 0{,}306$ et $P(\overline{M} \cap \overline{T}) = 0{,}30 \times 0{,}95 = 0{,}285$, on obtient :
    \[
    P_{\overline{T}}(\overline{M}) = \frac{0{,}285}{0{,}306} = \frac{285}{306} = \frac{95}{102} \approx 0{,}93137\dots
    \]
    Au millième près :
    \[
    \mathbf{\mathrm{VPN} = P_{\overline{T}}(\overline{M}) \approx 0{,}931} \quad (\text{soit } 93{,}1\,\%)
    \]
\end{itemize}

\textbf{5. (b) Comparaison et interprétation clinique :}
On constate que :
\[
\mathrm{VPP} \approx 0{,}978 > \mathrm{VPN} \approx 0{,}931
\]
\begin{itemize}[label=$\bullet$]
    \item Le taux d'erreur sur un test positif est de $1 - \mathrm{VPP} = 1 - 0{,}978 = 0{,}022$, soit seulement \textbf{$2{,}2\,\%$} de faux positifs parmi les individus testés positifs.
    \item Le taux d'erreur sur un test négatif est de $1 - \mathrm{VPN} = 1 - 0{,}931 = 0{,}069$, soit environ \textbf{$6{,}9\,\%$} de faux négatifs parmi les individus testés négatifs.
\end{itemize}
\textbf{Interprétation :} Un résultat positif est donc plus fiable et plus probant qu'un résultat négatif. Ce phénomène s'explique par la très forte prévalence de la maladie dans la population de coyotes d'Oklahoma ($70\,\%$) : la forte proportion initiale d'animaux infectés renforce la certitude d'un diagnostic positif, mais augmente le risque relatif qu'un animal infecté échappe au test (faux négatif).
\end{correctionbox}

\begin{alertebox}[-- Influence de la prévalence sur la valeur prédictive]
En épidémiologie, la sensibilité ($P_M(T)$) et la spécificité ($P_{\overline{M}}(\overline{T})$) sont des caractéristiques intrinsèques du test biologique, mais la $\mathrm{VPP}$ et la $\mathrm{VPN}$ dépendent de manière cruciale de la \textbf{prévalence} $p = P(M)$ :
\[
\mathrm{VPP}(p) = \frac{p \times \mathrm{Se}}{p \times \mathrm{Se} + (1-p)(1-\mathrm{Sp})}
\]
Dans une population où la maladie est rare ($p \to 0$), la $\mathrm{VPP}$ s'effondre vers $0$ même avec une excellente spécificité de $95\,\%$ (c'est le fameux \emph{paradoxe des faux positifs}). Ici, au contraire, la très forte prévalence ($p = 0{,}70$) tire la $\mathrm{VPP}$ à plus de $97{,}8\,\%$.
\end{alertebox}

\vspace{0.4cm}

% ----------------------------------------------------
% Correction Exercice 4
% ----------------------------------------------------
\subsection*{Correction détaillée de l'Exercice 4 \hfill (5 points)}

\begin{correctionbox}[-- Variable aléatoire, espérance, signe du trinôme \& Optimisation]
\textbf{1. Cas particulier : $2$ jetons noirs et $3$ jetons blancs ($N = 2$, total $5$ jetons)}

\textbf{(a) Arbre pondéré :}
Le tirage s'effectuant avec remise, les deux tirages sont indépendants et les probabilités restent invariantes :
\[
P(B) = \frac{3}{5} = 0{,}6 \quad \text{et} \quad P(N) = \frac{2}{5} = 0{,}4
\]
\begin{center}
\begin{tikzpicture}[xscale=2.6, yscale=1.1, thick, >=Stealth]
    \node[circle, fill=llgblue!20, inner sep=2pt] (R) at (0, 0) {$\Omega$};
    \node (B1) at (1.5, 0.8) {$B_1$};
    \node (N1) at (1.5, -0.8) {$N_1$};
    \node (B2B) at (3.3, 1.2) {$B_2$};
    \node (N2B) at (3.3, 0.4) {$N_2$};
    \node (B2N) at (3.3, -0.4) {$B_2$};
    \node (N2N) at (3.3, -1.2) {$N_2$};
    
    \draw[->] (R) -- (B1) node[midway, above left] {$3/5$};
    \draw[->] (R) -- (N1) node[midway, below left] {$2/5$};
    
    \draw[->] (B1) -- (B2B) node[midway, above] {$3/5$};
    \draw[->] (B1) -- (N2B) node[midway, below] {$2/5$};
    
    \draw[->] (N1) -- (B2N) node[midway, above] {$3/5$};
    \draw[->] (N1) -- (N2N) node[midway, below] {$2/5$};
\end{tikzpicture}
\end{center}

\textbf{(b) Probabilité de perdre $9$ \euro{} :}
Perdre $9$ euros correspond au tirage de deux jetons blancs consécutifs ($B_1 \cap B_2$) :
\[
P(X = -9) = P(B_1) \times P(B_2) = \left(\frac{3}{5}\right)^2 = \frac{9}{25} = \mathbf{0{,}36}
\]

\textbf{2. Cas général : $3$ jetons blancs et $N$ jetons noirs ($N \ge 2$, total $N+3$ jetons)}

\textbf{(a) Loi de probabilité de la variable aléatoire $X$ :}
Le gain algébrique $X$ prend ses valeurs dans l'ensemble $\{-9, -1, 5\}$.
Les probabilités associées s'expriment en fonction de $N$ :
\begin{itemize}[label=$\bullet$]
    \item \textbf{Deux jetons blancs ($X = -9$) :}
    \[
    P(X = -9) = \left(\frac{3}{N+3}\right)^2 = \mathbf{\frac{9}{(N+3)^2}}
    \]
    \item \textbf{Deux jetons noirs ($X = -1$) :}
    \[
    P(X = -1) = \left(\frac{N}{N+3}\right)^2 = \mathbf{\frac{N^2}{(N+3)^2}}
    \]
    \item \textbf{Deux jetons de couleurs différentes ($X = 5$) :}
    Il y a deux ordres possibles : $(B_1 \cap N_2)$ ou $(N_1 \cap B_2)$ :
    \[
    P(X = 5) = 2 \times \frac{3}{N+3} \times \frac{N}{N+3} = \mathbf{\frac{6N}{(N+3)^2}}
    \]
\end{itemize}
\emph{Vérification de la somme :}
\[
\sum_{x \in X(\Omega)} P(X = x) = \frac{9 + N^2 + 6N}{(N+3)^2} = \frac{(N+3)^2}{(N+3)^2} = 1
\]
La loi de probabilité de $X$ est résumée dans le tableau suivant :
\begin{center}
\renewcommand{\arraystretch}{1.5}
\begin{tabular}{|c|c|c|c|}
    \hline
    $x_i$ & $-9$ & $-1$ & $5$ \\
    \hline
    $P(X = x_i)$ & $\dfrac{9}{(N+3)^2}$ & $\dfrac{N^2}{(N+3)^2}$ & $\dfrac{6N}{(N+3)^2}$ \\
    \hline
\end{tabular}
\end{center}

\textbf{(b) Résolution de l'inéquation $-x^2 + 30x - 81 > 0$ :}
Considérons le trinôme du second degré $P(x) = -x^2 + 30x - 81$.
Son discriminant vaut :
\[
\Delta = b^2 - 4ac = 30^2 - 4(-1)(-81) = 900 - 324 = 576 = 24^2
\]
Comme $\Delta > 0$, le trinôme admet deux racines réelles distinctes :
\[
x_1 = \frac{-b - \sqrt{\Delta}}{2a} = \frac{-30 - 24}{-2} = \frac{-54}{-2} = 27
\]
\[
x_2 = \frac{-b + \sqrt{\Delta}}{2a} = \frac{-30 + 24}{-2} = \frac{-6}{-2} = 3
\]
Le coefficient dominant étant $a = -1 < 0$, la parabole est orientée vers le bas. Le trinôme est donc strictement positif à l'intérieur de ses racines :
\[
-x^2 + 30x - 81 > 0 \iff \mathbf{x \in \,]3, 27[}
\]

\textbf{(c) Nombre de jetons noirs pour un jeu favorable :}
Le jeu est favorable au joueur si et seulement si son espérance mathématique est strictement positive : $\E(X) > 0$.
Calculons $\E(X)$ :
\[
\E(X) = \sum_{i=1}^3 x_i P(X = x_i) = (-9) \times \frac{9}{(N+3)^2} + (-1) \times \frac{N^2}{(N+3)^2} + 5 \times \frac{6N}{(N+3)^2}
\]
\[
\E(X) = \frac{-81 - N^2 + 30N}{(N+3)^2} = \mathbf{\frac{-N^2 + 30N - 81}{(N+3)^2}}
\]
Comme $N \ge 2$, le dénominateur $(N+3)^2$ est strictement positif. Par conséquent :
\[
\E(X) > 0 \iff -N^2 + 30N - 81 > 0
\]
D'après la question 2.(b), cette condition équivaut à $N \in \,]3, 27[$.
Puisque le nombre de jetons noirs $N$ est un entier naturel, le jeu est favorable au joueur si et seulement si :
\[
N \in \boldsymbol{[\![4, 26]\!]}
\]
Il faut donc que l'urne contienne entre \textbf{$4$ et $26$ jetons noirs inclus}.

\textbf{(d) Nombre de jetons noirs maximisant le gain moyen :}
Le gain moyen est modélisé par la fonction $f$ définie sur $[2, +\infty[$ par :
\[
f(N) = \frac{-N^2 + 30N - 81}{(N+3)^2}
\]
La fonction $f$ est rationnelle, continue et dérivable sur $[2, +\infty[$. Posons :
\[
u(N) = -N^2 + 30N - 81 \implies u'(N) = -2N + 30
\]
\[
v(N) = (N+3)^2 \implies v'(N) = 2(N+3)
\]
Calculons la dérivée $f'(N)$ en factorisant par $(N+3)$ au numérateur :
\[
f'(N) = \frac{u'(N)v(N) - u(N)v'(N)}{(v(N))^2} = \frac{(-2N + 30)(N+3)^2 - 2(N+3)(-N^2 + 30N - 81)}{(N+3)^4}
\]
\[
f'(N) = \frac{(N+3)\left[(-2N + 30)(N+3) - 2(-N^2 + 30N - 81)\right]}{(N+3)^4}
\]
Développons l'expression entre crochets :
\[
(-2N^2 - 6N + 30N + 90) - (-2N^2 + 60N - 162) = -2N^2 + 24N + 90 + 2N^2 - 60N + 162 = -36N + 252
\]
D'où :
\[
f'(N) = \frac{-36N + 252}{(N+3)^3} = \mathbf{\frac{-36(N - 7)}{(N+3)^3}}
\]
Étudions le signe de $f'(N)$ sur $[2, +\infty[$ :
\begin{itemize}[label=$\bullet$]
    \item Pour $N \in [2, 7[$, $-36(N-7) > 0$ donc $f'(N) > 0$ : la fonction $f$ est strictement croissante.
    \item En $N = 7$, $f'(7) = 0$.
    \item Pour $N > 7$, $-36(N-7) < 0$ donc $f'(N) < 0$ : la fonction $f$ est strictement décroissante.
\end{itemize}
La fonction $f$ admet donc un maximum absolu strict en $N = 7$.
L'entier recherché est donc :
\[
\mathbf{N = 7}
\]
Le joueur doit exiger exactement \textbf{$7$ jetons noirs} dans l'urne.
Le gain moyen maximal réalisable s'élève à :
\[
\E_{\max} = f(7) = \frac{-(7)^2 + 30(7) - 81}{(7+3)^2} = \frac{-49 + 210 - 81}{100} = \frac{80}{100} = \mathbf{0{,}80\text{ \euro}}
\]
\end{correctionbox}

\vspace{0.4cm}

% ----------------------------------------------------
% Correction Exercice 5
% ----------------------------------------------------
\subsection*{Correction détaillée de l'Exercice 5 \hfill (1,5 point)}

\begin{correctionbox}[-- Problème de Mosteller : Fiabilité comparée de systèmes redondants]
Soit $p \in [0, 1]$ la probabilité élémentaire qu'un moteur tombe en panne. La fiabilité de chaque moteur est donc $q = 1-p$.
Les pannes étant mutuellement indépendantes, le nombre de pannes suit une loi binomiale.

Un avion arrive à destination si \textbf{strictement moins de la moitié} de ses moteurs tombent en panne :
\begin{itemize}[label=$\bullet$]
    \item \textbf{Avion A ($2$ moteurs) :} Strictement moins de la moitié de $2$ signifie $0$ panne.
    L'avion A arrive à bon port si et seulement si ses deux moteurs fonctionnent.
    Notons $X$ le nombre de moteurs en panne : $X \sim \mathcal{B}(2, p)$.
    \[
    P_A = P(X = 0) = (1-p)^2
    \]
    \item \textbf{Avion B ($4$ moteurs) :} Strictement moins de la moitié de $4$ signifie $0$ ou $1$ panne.
    Notons $Y$ le nombre de moteurs en panne : $Y \sim \mathcal{B}(4, p)$.
    \[
    P_B = P(Y = 0) + P(Y = 1) = \binom{4}{0} (1-p)^4 + \binom{4}{1} p(1-p)^3 = (1-p)^4 + 4p(1-p)^3
    \]
    En factorisant par $(1-p)^3$ :
    \[
    P_B = (1-p)^3 \left[(1-p) + 4p\right] = (1-p)^3 (1 + 3p)
    \]
\end{itemize}

\textbf{Comparaison des fiabilités $P_B - P_A$ :}
Pour déterminer quel avion offre la plus grande sécurité, étudions le signe de la différence $D(p) = P_B - P_A$ sur $[0, 1]$ :
\[
D(p) = (1-p)^3 (1+3p) - (1-p)^2 = (1-p)^2 \left[(1-p)(1+3p) - 1\right]
\]
Développons le crochet :
\[
(1-p)(1+3p) - 1 = 1 + 3p - p - 3p^2 - 1 = 2p - 3p^2 = p(2 - 3p)
\]
D'où la factorisation remarquable :
\[
\mathbf{D(p) = P_B - P_A = p (1-p)^2 (2 - 3p)}
\]

\textbf{Discussion exhaustive selon la valeur de $p$ :}
Puisque $p \ge 0$ et $(1-p)^2 \ge 0$ sur $[0, 1]$, le signe de $D(p)$ est exclusivement gouverné par le facteur affine $2 - 3p$ :
\begin{enumerate}
    \item \textbf{Cas d'indifférence ($D(p) = 0$) :}
    \[
    D(p) = 0 \iff p = 0 \quad \text{ou} \quad p = 1 \quad \text{ou} \quad p = \frac{2}{3}
    \]
    \begin{itemize}[label=$\bullet$]
        \item Si $p = 0$ (moteurs parfaits), $P_A = P_B = 1$.
        \item Si $p = 1$ (moteurs systématiquement défaillants), $P_A = P_B = 0$.
        \item Si $p = \frac{2}{3}$, $P_A = \left(\frac{1}{3}\right)^2 = \frac{1}{9}$ et $P_B = \left(\frac{1}{3}\right)^3 (1+2) = \frac{3}{27} = \frac{1}{9}$.
    \end{itemize}
    Dans ces trois cas, le choix de l'avion est \textbf{indifférent}.

    \item \textbf{Cas où les moteurs sont fiables ($p \in \left]0, \frac{2}{3}\right[$) :}
    Alors $2 - 3p > 0$, donc $D(p) > 0 \iff P_B > P_A$.
    \textbf{On choisit impérativement l'avion B (à $4$ moteurs)}.
    \emph{Exemple pratique :} Pour des moteurs aéronautiques modernes très fiables ($p = 0{,}01$), l'avion B offre une probabilité de crash de l'ordre de $6 \times 10^{-4}$ contre $2 \times 10^{-2}$ pour l'avion A.

    \item \textbf{Cas où les moteurs sont très défaillants ($p \in \left]\frac{2}{3}, 1\right[$) :}
    Alors $2 - 3p < 0$, donc $D(p) < 0 \iff P_A > P_B$.
    \textbf{On choisit l'avion A (à $2$ moteurs)}.
\end{enumerate}

\textbf{Conclusion synthétique :}
\[
\begin{cases}
    \text{Si } 0 < p < \dfrac{2}{3}, & \textbf{Choisir l'avion B (4 moteurs)} \\[0.25cm]
    \text{Si } p = 0, \; p = \dfrac{2}{3} \text{ ou } p = 1, & \textbf{Choix indifférent} \\[0.25cm]
    \text{Si } \dfrac{2}{3} < p < 1, & \textbf{Choisir l'avion A (2 moteurs)}
\end{cases}
\]
\end{correctionbox}

\begin{remarquebox}[-- Regard Prépa : Pourquoi plus de moteurs peut-il être néfaste ?]
Ce problème, introduit par le mathématicien et statisticien Frederick Mosteller dans \emph{Fifty Challenging Problems in Probability} (1965), illustre la dialectique fondamentale de la théorie de la fiabilité :
\begin{itemize}[label=$\bullet$]
    \item Multiplier les composants apporte de la \textbf{redondance passive} bénéfique lorsque le composant élémentaire est robuste ($p < 2/3$).
    \item Mais chaque moteur supplémentaire représente également une \textbf{source additionnelle de défaillance}. Si les moteurs sont intrinsèquement mauvais ($p > 2/3$), avoir 4 moteurs augmente considérablement le risque d'en voir tomber au moins 2 en panne ($P(Y \ge 2)$) par rapport au risque de voir tomber les 2 seuls moteurs d'un bi-moteur ($P(X = 2)$).
\end{itemize}
Le seuil critique $p^* = 2/3$ constitue le point de bifurcation où la redondance bascule de la protection vers la vulnérabilité.
\end{remarquebox}

\end{document}
